\section{Evaluation}
\label{sec:evaluation}

% PENDING RE-MEASUREMENT. The numbers presently in this section come from the
% held-out battery of 2026-07-29, which the 2026-08-06 audit invalidated: query
% radii were set by a spacing estimator that measured a sparsified cloud (up to
% 5x too large at 25M), the kd and grid tuning grids were depth-capped so their
% capacity parameter was inert, and every runner executed a binary frozen at
% 2026-07-28. All of it is being re-measured over the curated multi-scale
% benchmark. Do not quote a number from here until that lands; see testing.md.
% Reproduction after the re-run: scripts/make_cell_config.py, then
% scripts/prepare_curated_cells.py, then scripts/run_matrix_heatmap.ps1 and
% scripts/run_experiment_battery_v2.ps1 (both -CellConfig driven).

All experiments run on a single workstation (AMD-class multicore CPU, NVIDIA RTX 4080 SUPER, Windows 11); CPU evaluation is used wherever results are compared against CPU libraries. Workloads are replayed traces (Section~\ref{sec:trace-replay}); every candidate schema and every baseline answers the identical query sequence, and returned result counts are verified to match across all compared systems.

\paragraph{Protocol.} Three controls guard the comparisons. \emph{Held-out splits}: every recorded trace is partitioned into disjoint optimization and test halves, stratified by pipeline stage (for the batch regime, whole blocks are held out); all searching and tuning consumes the optimization half, and every number reported below is measured on the untouched test half with five repeat measurements of the full query batch. \emph{Tuned single-structure baselines}: ``best single'' never means a default configuration --- each single primitive receives a budget-matched grid over depth, leaf capacity, and axis policy drawn from the same design space the schema search uses, selected on the optimization half like any other candidate. \emph{Optimizer ablation}: the evolutionary search is compared against uniform random sampling of the same schema grammar under an identical number of full-fidelity evaluations (Section~\ref{sec:eval-ablation}).

\subsection{Pipeline replay: the advantage is scale-dependent}
\label{sec:eval-pipeline}

Table~\ref{tab:pipeline} reports the native-density pipeline cells. Two regularities organize the results. First, \emph{scale drives the searched advantage}: on the same terrain family, the searched schema is statistically indistinguishable from the best tuned single at 5M points, wins by 24\% (confidence intervals separated) at 50M, and the industrial 700M cell yields a $2.13\times$ separation in which the nine best full-scale candidates are all nested schemas. Second, \emph{the cost of a wrong fixed choice dwarfs the margin among good ones}: the worst tree-based default costs $4$--$13\times$ the best default on every cell (a regular grid, $148$--$777\times$), and the best default itself changes identity across cells --- the quadtree that wins terrain is $4\times$ off at Alhambra. The search never lost a cell: where a single structure is optimal, it returns one (at 5M it returned the plain quadtree itself), which is the ``allowed to choose a simple answer'' property of Section~\ref{sec:introduction} operating as designed.

\begin{table*}[t]
  \centering
  \caption{Pipeline-replay results on held-out queries (warm mean ms/query over five repeats). ``Best tuned single'' is the winner of a budget-matched per-primitive grid, not a default. CI-separated entries are bold.}
  \label{tab:pipeline}
  \small
  \begin{tabular}{llllll}
    \toprule
    dataset & searched winner & ms & best tuned single & ms & verdict \\
    \midrule
    SanAndreas 5M (terrain)       & quadtree (search returns single) & 0.0032 & qt$_{8}$/32 ($\equiv$ default) & 0.0032 & parity \\
    SanAndreas 50M (terrain)      & qt$_{8}$/256 + ot$_{2}$/128 + bvh$_{2}$/4k & \textbf{0.0081} & qt$_{8}$/32 & 0.0107 & \textbf{$-24\%$} \\
    Alhambra 100M (architecture)  & bvh$_{7}$/128 + ot$_{5}$/1k & 0.0130 & ot$_{8}$/1k & 0.0133 & $-2.5\%$ (n.s.) \\
    SolarPanels 700M (industrial) & qt$_{4}$/2k + ot$_{6}$/256 & \textbf{0.0377} & octree default$^{\ast}$ & 0.0803 & \textbf{$2.13\times$} \\
    \bottomrule
  \end{tabular}

  \smallskip
  \raggedright\footnotesize $^{\ast}$A full tuned-singles sweep at 700M is computationally infeasible (each candidate is a 700M-point build); the tuned sweep ran on a 5M subsample and its winner transferred poorly, so the strongest available single-structure competitor at full scale is reported. Subscripts are tree depths; the second number is leaf capacity.
\end{table*}

% Shape x size heatmap of held-out speedup over the best tuned single.
% Values are generated from results/eval_traces/matrix/heatmap_summary.csv and are
% identical to Table~\ref{tab:matrix}; regenerate both from that CSV if it changes.
% Colour: diverging blue<->orange on the manuscript's own accent pair, neutral grey
% at the 1.0 parity midpoint, so the parity boundary is the visual centre. Every
% cell also carries its printed value, so the encoding is never colour-alone
% (greyscale printing and CVD both stay legible).
\begin{figure*}[t]
  \centering
  \begin{tikzpicture}[x=17mm, y=8.4mm]
    % --- cell fills: blue = searched wins, orange = tuned single wins, grey = parity
    % \mdscell{col}{row}{value}{fill}
    \newcommand{\mdscell}[4]{%
      \fill[#4] (#1-0.5,#2-0.5) rectangle (#1+0.5,#2+0.5);
      \draw[white, line width=0.6pt] (#1-0.5,#2-0.5) rectangle (#1+0.5,#2+0.5);
      \node[font=\scriptsize, text=mdsink] at (#1,#2) {#3};}

    % row 5 = terrain
    \mdscell{1}{5}{0.88}{mdsout!42}
    \mdscell{2}{5}{0.96}{mdsout!16}
    \mdscell{3}{5}{1.09}{mdsaccent!24}
    \mdscell{4}{5}{1.27}{mdsaccent!62}
    % row 4 = industrial
    \mdscell{1}{4}{0.95}{mdsout!20}
    \mdscell{2}{4}{0.93}{mdsout!28}
    \mdscell{3}{4}{1.07}{mdsaccent!19}
    \mdscell{4}{4}{1.23}{mdsaccent!54}
    % row 3 = urban
    \mdscell{1}{3}{1.06}{mdsaccent!16}
    \mdscell{2}{3}{0.91}{mdsout!36}
    \mdscell{3}{3}{1.03}{mdsaccent!9}
    \mdscell{4}{3}{1.07}{mdsaccent!19}
    % row 2 = architecture
    \mdscell{1}{2}{0.99}{mdsfill}
    \mdscell{2}{2}{1.04}{mdsaccent!11}
    \mdscell{3}{2}{1.04}{mdsaccent!11}
    \node[font=\scriptsize, text=mdsline] at (4,2) {--};
    \draw[mdsline!40, line width=0.6pt] (3.5,1.5) rectangle (4.5,2.5);
    % row 1 = indoor
    \mdscell{1}{1}{0.96}{mdsout!16}
    \foreach \c in {2,3,4} {
      \node[font=\scriptsize, text=mdsline] at (\c,1) {--};
      \draw[mdsline!40, line width=0.6pt] (\c-0.5,0.5) rectangle (\c+0.5,1.5);}

    % --- axes
    \foreach \c/\lab in {1/1M, 2/5M, 3/25M, 4/100M}
      \node[font=\scriptsize, text=mdsink] at (\c,5.78) {\lab};
    \foreach \r/\lab in {5/terrain, 4/industrial, 3/urban (ALS), 2/architecture, 1/indoor}
      \node[font=\scriptsize, text=mdsink, anchor=east] at (0.42,\r) {\lab};

    % --- the crossover band: the single most important feature of the figure.
    % The rule stops below the size labels so it never collides with them; the
    % annotation sits clear above them.
    \draw[mdsaccent, thick, dashed] (2.5,0.42) -- (2.5,5.52);
    \node[font=\tiny\itshape, text=mdsaccent, anchor=south] at (2.5,6.15)
      {crossover: 10--25M};

    % --- legend
    \begin{scope}[shift={(0.5,-0.55)}]
      \fill[mdsout!42]     (0,0)     rectangle ++(0.30,0.30);
      \fill[mdsout!16]     (0.30,0)  rectangle ++(0.30,0.30);
      \fill[mdsfill]       (0.60,0)  rectangle ++(0.30,0.30);
      \fill[mdsaccent!24]  (0.90,0)  rectangle ++(0.30,0.30);
      \fill[mdsaccent!62]  (1.20,0)  rectangle ++(0.30,0.30);
      \draw[mdsline!50] (0,0) rectangle (1.50,0.30);
      \node[font=\tiny, text=mdsink, anchor=west] at (1.60,0.15)
        {$0.88\ \longleftarrow$ tuned single wins $\ \cdot\ 1.0\ \cdot\ $ searched wins $\longrightarrow\ 1.27$};
    \end{scope}
  \end{tikzpicture}
  \caption{Held-out speedup of the searched schema over the best tuned single, per
  matrix cell (the values of Table~\ref{tab:matrix}). Blue marks cells where the
  searched schema wins, orange where a tuned single does, and the printed value
  carries the magnitude so the encoding never rests on colour alone. The crossover
  from parity to searched wins falls between 10M and 25M points for every
  morphology, which is the scale law the evaluation rests on; shape modulates only
  the magnitude above it. Dashes mark cells with no source cloud at that size.}
  \label{fig:matrix-heatmap}
\end{figure*}


\subsection{Shape $\times$ size matrix}
\label{sec:eval-matrix}

The native-density cells confound shape with scale, so Table~\ref{tab:matrix} evaluates a controlled matrix: five source distributions subsampled to a common size ladder, each cell searched on its optimization half and scored on its test half. The crossover from parity to searched wins occurs between 10M and 25M points for \emph{every} shape, confirming scale as the first-order driver with shape held constant; shape then modulates the magnitude at 100M (terrain $1.27\times$, industrial $1.23\times$, urban $1.07\times$, architecture flat). The mechanism also varies: most winning schemas at scale are nested, but the terrain-100M winner is a \emph{single} quadtree at depth 9 --- a configuration between the grid points of the tuned-singles sweep --- showing that the search pays off through two routes, composition where depth regimes exist and off-grid parameter tuning where they do not. Because matrix cells are uniform subsamples, point count and density co-vary along each row; the native-density ladder of Table~\ref{tab:pipeline} exhibits the same trend, which addresses this confound.

\begin{table}[t]
  \centering
  \caption{Held-out speedup of the searched schema over the best tuned single per matrix cell (values $>1$: searched wins). Crossover occurs at 10--25M for every shape.}
  \label{tab:matrix}
  \small
  \begin{tabular}{lcccc}
    \toprule
    shape $\backslash$ size & 1M & 5M & 25M & 100M \\
    \midrule
    terrain      & 0.88 & 0.96 & \textbf{1.09} & \textbf{1.27} \\
    industrial   & 0.95 & 0.93 & \textbf{1.07} & \textbf{1.23} \\
    urban (ALS)  & 1.06 & 0.91 & 1.03 & \textbf{1.07} \\
    architecture & 0.99 & 1.04 & 1.04 & ---$^{\ast}$ \\
    indoor       & 0.96 & --- & --- & --- \\
    \bottomrule
  \end{tabular}

  \smallskip
  \raggedright\footnotesize $^{\ast}$The architecture source (Alhambra) falls 0.2\% short of the 100M rung; its full-cloud result is the Alhambra row of Table~\ref{tab:pipeline}. Indoor is limited to 1M by source size.
\end{table}

\subsection{Batch distribution (deep-learning dataloaders)}
\label{sec:eval-batch}

The batch regime rebuilds the index per block (32 PointNet++-style blocks per dataset, three point-set/trace pairs per block), scoring candidates by full-pass build-plus-query wall cost. Here the held-out protocol materially corrected an earlier conclusion: with training and evaluation on the \emph{same} blocks and default-configured baselines, searched schemas appeared to cut a full pass by 13--17\%; with 16 held-out blocks and tuned baselines, a tuned shallow octree (depth 4, leaf capacity 256) wins both datasets outright, with the searched schemas within 1--3\% of it (Table~\ref{tab:batch}). The blocks are simply too small ($\sim$4k points) for composition to matter --- consistent with the scale law of Table~\ref{tab:matrix} --- and the regime's honest story is \emph{selection plus insurance}: the search converges to near block-optimal without prior knowledge, while a wrongly fixed structure costs up to $1.9\times$ per pass, every batch, every epoch.

\begin{table*}[t]
  \centering
  \caption{Batch-distribution results on 16 held-out blocks (full-pass build + query wall cost, ms; lower is better).}
  \label{tab:batch}
  \small
  \begin{tabular}{lllll}
    \toprule
    dataset & best (tuned single) & pass & searched schema & worst default \\
    \midrule
    SanAndreas blocks & octree$_{4}$/256 & \textbf{116.3} & 118.7 ($+2.1\%$) & 223.4 ($1.9\times$) \\
    Alhambra blocks   & octree$_{4}$/256 & \textbf{115.0} & 116.5 ($+1.3\%$) & 241.2 ($1.9\times$) \\
    \bottomrule
  \end{tabular}
\end{table*}

\subsection{Optimizer ablation: sampling suffices}
\label{sec:eval-ablation}

Under budget-matched full-fidelity evaluations, uniform random sampling of the schema grammar matched or exceeded the evolutionary search (with and without diagnostic repair mutations) on the held-out half of every cell; at 50M the random arm's selection was the overall winner with separated confidence intervals, and the 700M full-scale confirmation is led by uniformly sampled candidates. We therefore attribute the system's gains to the \emph{schema grammar} (which spans single, nested, and conditional designs at all parameterizations) combined with \emph{measured selection on the real workload}, rather than to any particular traversal of that space. The evolutionary machinery is retained as a convenience, but the result is a simplification: practitioners can reproduce the wins with plain sampling plus measurement. Among near-tied candidates the selected winner varies with measurement noise across repeated searches, so we report winner \emph{families} rather than exact schemas in parity cells.

\subsection{Cross-cloud transfer}
\label{sec:eval-transfer}

Tuning is per-instance, but a natural question is what a tuned schema costs when applied to a different cloud of the same morphology. Replaying the SanAndreas-50M winner on the held-out queries of an independent 25M coastal-terrain survey (San Simeon, different sensor campaign and density) ranks it second at $+9\%$ over that cloud's best structure, while wrongly fixed singles cost $1.4$--$6.8\times$. Transferred schemas are thus \emph{safe but not optimal}: reusing a same-morphology schema is a reasonable default, and re-tuning recovers the last ten percent.

\subsection{External library baselines}
\label{sec:eval-frameworks}

Table~\ref{tab:frameworks} compares the searched schema against the indexes practitioners actually use, on identical held-out queries with verified result counts (residual disagreements are $\leq 1$ point on $\leq 10$ of 3000 queries, attributable to boundary rounding). The comparison exhibits the same crossover as everything above: at 5M, PCL's FLANN kd-tree is genuinely the fastest structure ($0.85\times$ our latency) --- small clouds are mature-library territory, and they are also exactly the parity cells where instance optimization has nothing to offer; at 100M the searched schema is $2.76\times$ faster than FLANN and $3.7\times$ faster than Open3D. FLANN builds $5.4\times$ faster at 100M (12.4\,s vs.\ 67.5\,s), so the crossover amortizes after $\sim$2.4M queries --- 2.4\% of a single per-point pipeline stage on that cloud. PDAL is omitted (box-crop-only harness support). As a primitive-quality control we also replayed identical range, radius and kNN queries against the bulk-loaded octree, kd-tree and packed R-tree of the Indexicon library~\cite{indexicon2026}. Exactness first: returned counts matched on every one of 351{,}000 sampled queries (13 cells $\times$ 3 structures $\times$ 3000 queries $\times$ 3 repeats), so the two systems demonstrably answer the same workload. On latency our primitives reach parity by 25M points. We previously reported them as faster on every comparable operation; that comparison is withdrawn, because it was measured against a stale binary, and because the per-node instrumentation described in Section~\ref{sec:disc-locality} was inside our timed region but never inside theirs. Indexicon's builds are $5$--$7\times$ faster, consistent with its focus on lean, portable, dynamic structures rather than query-optimized static ones. One result from this control is worth stating on its own: Indexicon's packed R-tree was the slowest of its three structures in 8 of the 13 cells and the fastest in none, which is external evidence for a grammar decision we had otherwise made on internal grounds --- overlapping minimum bounding rectangles do not fit a space of per-node disjoint splits, and we exclude the R-tree family accordingly.

\begin{table*}[t]
  \centering
  \caption{External baselines on identical held-out pipeline queries (mean ms/query; result counts verified). Bold marks the fastest per row.}
  \label{tab:frameworks}
  \small
  \begin{tabular}{lllll}
    \toprule
    dataset & searched schema & Open3D KDTreeFlann & PCL kd-tree & PCL octree \\
    \midrule
    SanAndreas 5M  & 0.0036 & 0.0133 ($3.7\times$) & \textbf{0.0028} ($0.85\times$) & 0.0147 \\
    Alhambra 100M  & \textbf{0.0129} & 0.0479 ($3.7\times$) & 0.0356 ($2.76\times$) & 0.0590 \\
    \bottomrule
  \end{tabular}
\end{table*}

% Amortization of the one-off build cost against query volume, Alhambra 100M cell.
% Source: results/framework_compare/v2_alhambra_100m_pcl/framework_comparison.md
%   ours  : build 67 453.68 ms, 0.012930 ms/query
%   FLANN : build 12 385.20 ms, 0.035645 ms/query
% Break-even = (67.4537 - 12.3852) / (3.5645e-5 - 1.2930e-5) = 2 424 322 queries.
% Curve coordinates are precomputed (scratchpad/gen_amort.py) because TeX overflows
% on 10^x. y is sqrt-compressed over 0..3600 s so both the early build gap and the
% late linear growth stay readable in one panel; curves clamp at the 60-min ceiling.
% Colours are the manuscript's accent pair, validated as a diverging/emphasis pair,
% and both lines are directly labelled so identity never rests on colour alone.
\begin{figure*}[t]
  \centering
  \begin{tikzpicture}[x=1.55mm, y=1.15mm]
    % gridlines behind the curves
    \foreach \yy/\lab in {16.330/10, 28.284/30, 40.000/60} {
      \draw[mdsline!22, very thin] (0,\yy) -- (62,\yy);
      \draw[mdsline] (0,\yy) -- (-1.2,\yy);
      \node[font=\tiny, text=mdsink, anchor=east] at (-1.5,\yy) {\lab\,min};}
    \node[font=\tiny, text=mdsink, anchor=east] at (-1.5,0) {0};

    % axes
    \draw[mdsline] (0,0) -- (62,0);
    \draw[mdsline] (0,0) -- (0,40);

    \foreach \e/\lab in {0/4, 12.4/5, 24.8/6, 37.2/7, 49.6/8, 62/9} {
      \draw[mdsline] (\e,0) -- (\e,-1.2);
      \node[font=\tiny, text=mdsink, anchor=north] at (\e,-1.4) {$10^{\lab}$};}
    \node[font=\scriptsize, text=mdsink, anchor=north] at (31,-6.5)
      {queries issued against the index};
    \node[font=\scriptsize, text=mdsink, rotate=90] at (-11,20)
      {cumulative wall cost};

    % one full pipeline pass over the 100M cloud = 3e8 queries. The annotation is
    % anchored inside the panel so it cannot overflow the right-hand column edge.
    \draw[mdsaccent!45, dotted, thick] (55.516,0) -- (55.516,30);
    \node[font=\tiny\itshape, text=mdsaccent, anchor=north west, align=left]
      at (56.4,15.5) {one pipeline pass\\($3{\times}10^{8}$ queries)};

    % FLANN kd-tree: cheaper to build, dearer per query
    \draw[mdsout, very thick, dashed]
      plot coordinates {(0.000,2.380) (3.100,2.405) (6.200,2.451) (9.300,2.529)
      (12.400,2.662) (15.500,2.885) (18.600,3.243) (21.700,3.796) (24.800,4.620)
      (27.900,5.803) (31.000,7.457) (34.100,9.726) (37.200,12.803) (40.300,16.948)
      (43.400,22.505) (46.500,29.940) (49.600,39.871) (49.900,40.000)};

    % searched schema: dearer to build, cheaper per query
    \draw[mdsaccent, very thick]
      plot coordinates {(0.000,5.481) (3.100,5.485) (6.200,5.492) (9.300,5.505)
      (12.400,5.528) (15.500,5.568) (18.600,5.639) (21.700,5.763) (24.800,5.977)
      (27.900,6.340) (31.000,6.939) (34.100,7.893) (37.200,9.351) (40.300,11.497)
      (43.400,14.550) (46.500,18.792) (49.600,24.590) (52.700,32.433)
      (55.300,40.000)};

    % break-even
    \draw[mdsink!55, dotted] (29.569,0) -- (29.569,6.627);
    \fill[mdsink] (29.569,6.627) circle (0.9pt);
    \node[font=\tiny\itshape, text=mdsink, anchor=south east] at (28.4,10.5)
      {break-even $\approx2.4$M};

    % direct labels
    \node[font=\tiny, text=mdsout, anchor=east] at (46.0,31.6) {PCL FLANN kd-tree};
    \node[font=\tiny, text=mdsaccent, anchor=west] at (50.6,22.0) {searched schema};
  \end{tikzpicture}
  \caption{Why the one-off cost is free in this regime, on the Alhambra 100M cell.
  The searched schema builds more slowly (67.5\,s vs.\ 12.4\,s) but answers each
  query $2.76\times$ faster, so it repays the difference after roughly 2.4 million
  queries. One three-stage pipeline pass over that cloud issues $3\times10^{8}$
  queries --- two orders of magnitude past break-even --- which is why the build
  gap is 2.4\% of a single pass. Below the crossover the ordering reverses: at 5M
  the mature library is faster per query, and no query volume repays the
  difference.}
  \label{fig:amortization}
\end{figure*}


\subsection{Search cost}
\label{sec:eval-search-cost}

Table~\ref{tab:searchcost} quantifies the mitigations of Section~\ref{sec:app-search-cost}. The serial Alhambra search (43 candidate builds at 100M) anchors the cost; eight-way candidate dispatch reduces a representative 5M search from 9.4 to 1.5 minutes; and multi-fidelity search on a 2--5M subsample ranks candidates with Spearman $\rho = 0.88$--$0.95$ against full scale, with the true full-scale winner inside the subsample top-3 in every measured case --- so the practical recipe is minutes of subsample search plus three confirmation builds, at zero measured regret. The 700M cell is only reachable this way: a 0.9-minute subsample search plus eleven full-scale confirmation builds (3.5\,h, dominated by index construction) produced the $2.13\times$ winner of Table~\ref{tab:pipeline}; an exhaustive full-scale search at that size is out of reach entirely.

\begin{table*}[t]
  \centering
  \caption{Search-cost mitigations (held-out winners; regret measured against the full-scale search result).}
  \label{tab:searchcost}
  \small
  \begin{tabular}{llll}
    \toprule
    configuration & wall time & regret & notes \\
    \midrule
    full-scale GA, serial (Alhambra 100M)   & 42.9 min & --- & anchor; $>$95\% is candidate builds \\
    5M search, \texttt{--parallel 8}        & 1.5 min (from 9.4) & within noise & winner family unchanged \\
    subsample search + top-3 confirm        & minutes + 3 builds & 0\% & $\rho=0.88$ (2M) / $0.95$ (5M) \\
    700M via subsample + top-11 confirm     & 0.9 min + 3.5 h & --- & only feasible route at this size \\
    \bottomrule
  \end{tabular}
\end{table*}

\subsection{Boundary study: ray tracing}
\label{sec:eval-boundary}

Table~\ref{tab:boundary} summarizes the triangle-mesh study of Section~\ref{sec:app-boundary}: top-level structures (uniform grid, octree, kd-tree, and nested stacks) wrapped around per-cell state-of-the-art compressed wide BVHs, evaluated by GPU path-tracing throughput with exactness validated against the single global BVH (identical hits and ray paths). On a synthetic axis-aligned scene a searched grid nearly doubles coherent-ray throughput; on a real scanned model the single global SAH BVH dominates every wrapper and every nested schema, and the optimizer converges to it. Where a hyper-engineered incumbent exists, the honest output of a measured search is the incumbent.

\begin{table*}[t]
  \centering
  \caption{Boundary study: GPU ray throughput (Mrays/s, higher is better) of searched top-level structures over per-cell compressed wide BVH leaves, vs.\ a single global BVH.}
  \label{tab:boundary}
  \small
  \begin{tabular}{lllll}
    \toprule
    scene & single global & best grid & best octree/kd & best nested \\
    \midrule
    synthetic box (coherent rays) & 13\,876 & \textbf{25\,610} (+85\%) & 12\,019 / 10\,748 & 13\,269 \\
    scanned dragon (coherent rays) & \textbf{3\,132} & 2\,946 ($-6\%$) & 2\,544 / 1\,359 & 1\,564 \\
    \bottomrule
  \end{tabular}
\end{table*}

\subsection{Threats to validity}
\label{sec:eval-threats}

Single-machine measurements; derived (not library-instrumented) pipeline traces, exact for self-query stages by construction; one sampling seed per search arm, with winner-family stability across repeated searches reported in place of exact-schema stability; the batch regime models the dataloader's neighborhood queries but not the surrounding training loop. Point count and density co-vary along each size ladder, because a tier is a uniform subsample of its scene rather than a scan at that density; the per-scene ladders are nested by construction, so a scaling curve varies size alone and not the scene, but not size alone and density together. Tiers above a scene's native count would contain interpolated points and are excluded from every reported cell, which caps several scenes below the top of the ladder. The largest industrial cell lacks a full-scale tuned-singles sweep for cost reasons. Coordinates are single-precision offsets from a per-cloud origin, leaving a floor of a few millimetres on scenes spanning tens of kilometres, which is reported alongside any exactness comparison on those scenes. Held-out splits are drawn from the same cloud and trace distribution; the cross-cloud experiment (Section~\ref{sec:eval-transfer}) is the only distribution-shift test. A coordinate-frame audit is part of the protocol: replayed traces are checked for non-degenerate result counts (a frame mismatch between capture and index manifests as universally empty radius queries), the replay path performs an automatic world-to-local transform with a bounds-coverage guard, and all cross-library comparisons verify returned counts --- a discipline that caught two frame-convention bugs in our own harness during this work.
